\documentclass[12pt,a4paper]{article}
%\usepackage{parskip} % Removes indent, adds spacing between paragraphs
%\usepackage[margin=2cm]{geometry}
%\title{Appropriate Title for This Text}
%\author{}
%\date{}

\begin{document}
%\maketitle
This document illustrates the basic structure of a simple \LaTeX{} article. 
An article will typically also carry a title, an author name, and a date. Rather than formatting these by hand, \LaTeX{} provides dedicated commands for the job: \texttt{title}, \texttt{author}, and \texttt{date} store this information in the preamble, and a single \texttt{maketitle} command—placed at the top of the document body prints the title block automatically. This separation of content from formatting is a recurring theme in \LaTeX{}, and it is one of the main reasons documents remain so easy to maintain.
The functionality of \LaTeX{} can be extended almost without limit through packages, which are loaded in the preamble with \texttt{usepackage}. Thousands of packages are available on CTAN, covering everything from mathematics and bibliographies to graphics, tables, and typography refinements.
For instance, although the default margins of the article class are quite generous, you can adjust them easily with the popular geometry package: The writing in the preamble gives you one-inch margins on all sides, and far more sophisticated layouts are possible with just a few additional options.

One of the most useful things to understand early on is how \LaTeX{} interprets whitespace in your source file. An empty line tells \LaTeX{} to start a new paragraph. The text that follows the break will be typeset as a fresh paragraph, typically with a small indent (or extra vertical space, depending on the document class) marking the transition.
Line breaks on their own, however, are treated entirely differently. If you continue typing on a new line without leaving a blank line in between, \LaTeX{} will simply keep everything in the same paragraph. You can see this principle at work in the preceding example:
several sentences were typed on separate lines of the source file, yet in the typeset
output they form one continuous paragraph.
This is a great convenience rather: it means you are free to break lines wherever you like in your editor, keeping the source readable, without ever affecting the final layout. As the previous demonstration shows, the content of a single paragraph may be divided across many lines of source text and still be rendered as one cohesive block.
Just as line breaks are ignored, so is excess horizontal spacing. \LaTeX{} automatically formats the text to fit the margins of the page, stretching and shrinking the interword space subtly to produce perfectly justified, evenly spaced lines. Consequently, adding        extra         spaces      between         words
accomplishes      nothing:         no        matter        how         many      times       you       press  the space bar, the output spacing between words remains identical. The spacing between words in this very sentence stays perfectly uniform, even though the source file contains unusual and wildly inconsistent spacing throughout.

When you want deliberate, meaningful spacing for example between a number and its unit, or after a sentence ending in a capital letter you should express it explicitly, for instance with the $\sim$ character for a non-breaking space or with commands such as \textbackslash, and \textbackslash; for thin and thick spaces.

Mathematics can be woven into running text just as naturally as words are. Inline expressions are enclosed in dollar signs $...$, and \LaTeX{} takes care of all the typographic subtleties that make math look right properly sized symbols, correct spacing around operators, and italicized variables. Writing mathematical symbols in the text, like the following expression for the area of a circle $\alpha = \pi r^2$, or another expression for the volume of a sphere, $V = \frac{4}{3} \pi r^3$ can be done in a very straightforward manner.
 Notice how the source reads almost exactly like what you want to see: \textbackslash \texttt{pi} produces $\pi$ , a caret $\hat{}$ raises a superscript, and $\frac{a}{b}$ builds a stacked fraction. For anything more complex, displayed equations between \[ ax^2 + bx + c = 0 \rightarrow \quad x = \frac{-b \pm \sqrt{b^2 - 4ac}}{2a} \] or inside an equation environment will set the math apart from the text on its own centered line, numbered automatically for easy reference.
\end{document}